% TOP_PROBLEM: {"id": 30000362, "title": "Deformations and Massless Spectra of Smooth Deligne–Mumford Stacks", "category_id": 16, "category": {"id": 16, "name": "physics", "display_name": "Mathematical Physics", "description": "Problems at the intersection of mathematics and physics.", "slug": "physics", "order_index": 16, "created_at": "2026-07-31T15:26:25.671Z"}, "status": "partially_solved", "statusQualification": "The original registry excerpt also includes infrared agreement of different quotient presentations, separately catalogued at ID 30000363. Both aspects of the excerpt are retained by this note. Ordinary stack deformation theory can fail to match the physical deformations, as the report's C^2/Z_2 example demonstrates."}
% FIRST_POSTED: 2026-10-11
% ENTRY_KIND: statement_only
% SUBMISSION: {"kind": "statement_only", "category": "top_problems", "problem_id": "30000362", "proposal_type": "established", "source": "Mini-Workshop: Heterotic Strings, Derived Categories, and Stacks (2005), Oberwolfach Reports, DOI 10.4171/owr/2005/53, pp. 3019--3060; official report source and mathematical context. https://doi.org/10.4171/owr/2005/53; Official EMS Press PDF of the same Oberwolfach report; PDF pages 18--20, printed pages 3036--3038. https://ems.press/content/serial-article-files/46027; GLSMs for gerbes and stacks; reference in the saved dated catalogue literature review. https://arxiv.org/abs/hep-th/0502053; derived categories and stacks in physics; reference in the saved dated catalogue literature review. https://doi.org/10.48550/arXiv.hep-th/0608056", "status_evidence": "Saved source status is partially_solved; retained unchanged. The official historical OWR source was inspected on 2026-10-11. The saved literature assessment is dated 2026-08-21 and is not a fresh comprehensive status verification. ", "status_checked_on": "2026-10-11", "duplicate_check": "Checked canonical identity and normalized formulations against the complete saved UnsolvedMath registry; related or repeated records: 30000363. Mathematical equivalence and current literature were not exhaustively searched; no new ID allocated.", "attribution": "Eric Sharpe, joint work with T. Pantev. Statement preparation and source audit assisted by OpenAI GPT-6 (Codex); no human mathematical authorship or independent proof verification claimed."}
% SOURCE_STATEMENT: "For string compactification on a smooth Deligne–Mumford stack, how should deformations and massless spectra be defined so that they agree with the physical theory?"
% Preparation and final self-audit: CONTRIBUTING.md prompts C and D; source/audit artifacts retained in the local physics_statements_2026_10_11 batch.
% Imported from: physics_statements.tex; UnsolvedMath ID 30000362
% !TeX program = lualatex
\documentclass[11pt]{article}
\usepackage[margin=25mm]{geometry}
\usepackage{fontspec}
\setmainfont{Latin Modern Roman}
\usepackage{amsmath,amssymb,mathtools}
\usepackage{unicode-math}
\setmathfont{Latin Modern Math}
\usepackage{xurl,hyperref}
\hypersetup{colorlinks=true,urlcolor=blue}
\setcounter{secnumdepth}{0}
\setlength{\parindent}{0pt}
\setlength{\parskip}{5pt}
\setlength{\emergencystretch}{3em}
\newcommand{\sref}[2]{\hyperlink{source:#1:#2}{[S#2]}}
\newcommand{\cataloguescope}[1]{\paragraph{Recorded target.}#1}
\newcommand{\R}{\mathbb R}
\newcommand{\C}{\mathbb C}
\newcommand{\N}{\mathbb N}
\newcommand{\Z}{\mathbb Z}
\newcommand{\Q}{\mathbb Q}
\newcommand{\F}{\mathbb F}
\newcommand{\eps}{\varepsilon}
\begin{document}
% UnsolvedMath ID 30000362; source catalogue code OWR-1114-003
\setcounter{section}{30000361}
\section{Deformations and Massless Spectra of Smooth Deligne–Mumford Stacks}\label{30000362}
{\small UnsolvedMath ID 30000362 · Mathematical Physics}
\subsection{Definitions and mathematical statement}

\paragraph{Shared notation.}$\mathbb N=\{1,2,\ldots\}$; the standard number fields have their usual meanings. All additional conventions are stated locally.


\paragraph{Definitions and conventions.} A smooth Deligne--Mumford stack is a geometric object locally presented by smooth spaces with finite stabilizer groups. A quotient presentation $[X/G]$ specifies a space and a group action; distinct presentations may describe the same stack. The report associates a gauged sigma model to physically realizable presentations. A deformation changes the background or couplings infinitesimally; a massless spectrum consists of zero-mass physical excitations. The inertia stack records pairs of a point and an automorphism of that point and enters the proposed cohomological spectrum calculation.

\paragraph{Source scope and qualifications.} The original registry excerpt also includes infrared agreement of different quotient presentations, separately catalogued at ID 30000363. Both aspects of the excerpt are retained by this note. Ordinary stack deformation theory can fail to match the physical deformations, as the report's C^2/Z\_2 example demonstrates.

\paragraph{Bundled original excerpt.} The supplied original excerpt also asks whether different quotient presentations of the same stack flow to the same physical infrared theory. That presentation-independence question is recorded separately as problem 30000363; no claim that it follows from a spectrum prescription is made.

\paragraph{Catalogue mathematical target.}
\paragraph{Catalogue wording retained for traceability.}
For string compactification on a smooth Deligne–Mumford stack, how should deformations and massless spectra be defined so that they agree with the physical theory?

\paragraph{Saved literature context.} The supplied catalogue's assessment dated 2026-08-21 records: Stack presentations have been incorporated into gauged linear sigma models and gerbe/string constructions, producing concrete conformal theories in many examples; a general classification of strings and conformal field theories on stacks remains open. The assessment is retained as dated provenance; it is not a new comprehensive literature or open-status certification.
\paragraph{Scope and source qualification.}The original registry excerpt also includes infrared agreement of different quotient presentations, separately catalogued at ID 30000363. Both aspects of the excerpt are retained by this note. Ordinary stack deformation theory can fail to match the physical deformations, as the report's C^2/Z_2 example demonstrates.
\subsection{Short English statement}
Define deformations and massless spectra for string theory on stacks so that they agree with the associated physical theories.
\subsection{Sources}
\begin{itemize}
\item[\hypertarget{source:30000362:1}{S1}] ULAM.AI and the UnsolvedMath Contributors, UnsolvedMath, record 30000362 (OWR-1114-003). Catalogue curation under CC BY 4.0; definitions and English exposition added locally.
\url{https://huggingface.co/datasets/ulamai/UnsolvedMath}.
\item[\hypertarget{source:30000362:2}{S2}] Mini-Workshop: Heterotic Strings, Derived Categories, and Stacks (2005), Oberwolfach Reports, DOI 10.4171/owr/2005/53, pp. 3019--3060; official report source and mathematical context.
\url{https://doi.org/10.4171/owr/2005/53}.
\item[\hypertarget{source:30000362:3}{S3}] Official EMS Press PDF of the same Oberwolfach report; PDF pages 18--20, printed pages 3036--3038.
\url{https://ems.press/content/serial-article-files/46027}.
\item[\hypertarget{source:30000362:4}{S4}] GLSMs for gerbes and stacks; reference in the saved dated catalogue literature review.
\url{https://arxiv.org/abs/hep-th/0502053}.
\item[\hypertarget{source:30000362:5}{S5}] derived categories and stacks in physics; reference in the saved dated catalogue literature review.
\url{https://doi.org/10.48550/arXiv.hep-th/0608056}.
\end{itemize}
\label{end:30000362}
\end{document}
